%%%%%%%%%%%%%%%%%%%%%%%%%%%%%%%%%%%%%%%%%%%%%%%%%%%%%%%%%%%%%%%%%%%%%%
% Lớp 10 | Mệnh đề. Tập hợp | Mệnh đề
% Dạng 0D1?1-3: Phủ định của một mệnh đề
% Số câu: 66  (sinh tự động bởi tach_cau.py; từ nay thầy sửa trực tiếp file này)
%%%%%%%%%%%%%%%%%%%%%%%%%%%%%%%%%%%%%%%%%%%%%%%%%%%%%%%%%%%%%%%%%%%%%%

%%%%%%%%%%%%%%%% TRẮC NGHIỆM 4 PHƯƠNG ÁN %%%%%%%%%%%%%%%%

% nguồn: MaHoaNH_TracNghiem_Lop10
\begin{ex}%[0D1N1-3]%[ma:8396dabd]
	Cho mệnh đề $P\colon$ \lq\lq Nếu $a$ và $b$ chia hết cho 3 thì $a+b$ cũng chia hết cho 3\rq\rq. Mệnh  đề nào là mệnh đề đảo của mệnh đề $P$?
	\choice {Nếu $a+b$ chia hết cho 3 thì $a$ chia hết cho 3}{Nếu  $a+b$ chia hết cho 3 thì $b$ chia hết cho 3}{Nếu $a$ và $b$ không chia hết cho 3 thì $a+b$ cũng không chia hết cho 3}{\True Nếu $a+b$ chia hết cho 3 thì $a$ và $b$ chia hết cho 3}
	 \loigiai{
		 Mệnh đề $P$ có dạng \lq\lq Nếu $A$ thì $B$\rq\rq. Mệnh đề đảo là \lq\lq Nếu $B$ thì $A$\rq\rq.
	 }
\end{ex}

% nguồn: MaHoaNH_TracNghiem_Lop10
\begin{ex}%[0D1N1-3]%[ma:9d1b2270]
	Cho mệnh đề \lq\lq$\forall  x\in\mathbb{R}$ $x^{2}-x+7<0$\rq\rq. Hỏi mệnh đề nào là mệnh đề phủ định của mệnh đề trên?
	\choice {\True $\exists  x\in\mathbb{R},x^{2}-x+7\ge0$}{ $\forall x\in\mathbb{R},x^{2}-x+7>0$}{ $\forall x\in\mathbb{R},x^{2}-x+7<0$}{ $\not\exists x\in\mathbb{R},x^{2}-x+7<0$}
	 \loigiai{
		Mệnh đề phủ định của mệnh đề trên là $\exists  x\in\mathbb{R}, x^{2}-x+7\ge0$.
	} 
\end{ex}

% nguồn: MaHoaNH_TracNghiem_Lop10
\begin{ex}%[0D1N1-3]%[ma:0721d4e0]
	Cho mệnh đề $P\colon "\forall x\in\mathbb{R},  x^{2}+1\ge2x"$. Mệnh đề phủ định $\bar{P}$ là
	\choice {$\forall x\in\mathbb{R}, x^{2}-2x+1<0$}{\True $\exists x\in\mathbb{R}, x^{2}-2x+1<0$}{$\exists  x\in\mathbb{R}, x^{2}+1\le2x$}{$\exists x\in\mathbb{R}, x^{2}+1\ge2x$}
	 \loigiai{
		 Mệnh đề phủ định của $P$ là $\bar{P}$ là $\exists x\in\mathbb{R}, x^{2}-2x+1<0$.
	 }
\end{ex}

% nguồn: MaHoaNH_TracNghiem_Lop10
\begin{ex}%[0D1N1-3]%[ma:1c6353d4]
	Mệnh đề nào sau đây có mệnh đề đảo đúng? 
	 \choice
		{Nếu $a=b$ thì $a^2=b^2$} 
		{\True Nếu một phương trình bậc hai có $\Delta > 0$ thì phương trình đó có $2$  nghiệm phân biệt}
		{Nếu một số chia hết cho $6$ thì cũng  chia hết cho $3$}
		{Nếu hai tam  giác bằng nhau thì có diện tích bằng nhau}
	 \loigiai{
		Nếu một phương trình bậc  hai có $\Delta > 0$ thì phương trình đó có $2$ nghiệm phân biệt có mệnh đề đảo đúng.
	} 
\end{ex}

% nguồn: MaHoaNH_TracNghiem_Lop10
\begin{ex}%[0D1N1-3]%[ma:6af0c2ff]
	Mệnh đề phủ định của mệnh đề \lq\lq Có ít nhất  một số thực $x$ mà bình phương của nó là 1 số âm\rq\rq\ là
	\choice 
		{\True \lq\lq $\forall x\in \mathbb{R}:x^2\ge 0$\rq\rq} 
		{\lq\lq $\exists x\in \mathbb{R}:x^2>  0$ \rq\rq}
		{\lq\lq $\forall x\in \mathbb{R}:x^2>  0$\rq\rq}
		 {\lq\lq $\exists x\in \mathbb{R}:x^2\ge 0$\rq\rq}
	\loigiai{ 
		Ta có: $P:$ \lq\lq $\exists x\in \mathbb{R}:x^2< 0$\rq\rq\, $\Rightarrow \overline{P}:$ \lq\lq $\forall x\in \mathbb{R}:x^2\ge 0$\rq\rq. 
	} 
\end{ex}

% nguồn: MaHoaNH_TracNghiem_Lop10
\begin{ex}%[0D1H1-3]%[ma:c1dbb495]
	Cho mệnh đề $P\colon$ \lq\lq$\exists x\in \mathbb{R}$, $x<\dfrac{1}{x}$\rq\rq. Xác định mệnh đề phủ  định của mệnh đề $P$.
	\choice 
		{$\overline{P}\colon$ \lq\lq $\exists x\in  \mathbb{R}$, $x\ge \dfrac{1}{x}$\rq\rq}
		{$\overline{P}\colon$ \lq\lq $\forall x\in  \mathbb{R}$, $x> \dfrac{1}{x}$\rq\rq}
		{\True $\overline{P}\colon$ \lq\lq $\forall  x\in \mathbb{R}$, $x\ge \dfrac{1}{x}$\rq\rq}
		 {$\overline{P}\colon$ \lq\lq $\exists x\in \mathbb{R}$, $x> \dfrac{1}{x}$\rq\rq}
	 \loigiai{
	Phủ  định của mệnh đề $P\colon$ \lq\lq$\exists x\in \mathbb{R}$, $x<\dfrac{1}{x}$\rq\rq\, là mệnh đề $\overline{P}\colon$ \lq\lq $\forall x\in \mathbb{R}$, $x\ge \dfrac{1}{x}$\rq\rq.}
\end{ex}

% nguồn: MaHoaNH_TracNghiem_Lop10
\begin{ex}%[0D1H1-3]%[ma:d4ff0c9f]
	Phủ định của mệnh đề \lq\lq Hình thoi có  hai đường chéo vuông góc với nhau\rq\rq\ là
	 \choice
		{\lq\lq Hai  đường chéo của hình thoi vuông góc với nhau\rq\rq}
		{\True\lq\lq Hình thoi  có hai đường chéo không vuông góc với nhau\rq\rq}
		{\lq\lq Hình  thoi có hai đường chéo bằng nhau\rq\rq}
		{\lq\lq Hình thoi là hình bình hành có hai đường  chéo vuông góc với nhau\rq\rq}
	\loigiai{ 
		Phủ định của \lq\lq vuông góc\rq\rq\ là \lq\lq không vuông  góc\rq\rq.
	 }
\end{ex}

% nguồn: MaHoaNH_TracNghiem_Lop10
\begin{ex}%[0D1H1-3]%[ma:c55eb03a]
	Mệnh đề phủ định của mệnh đề \lq\lq  $2018$ là số nguyên tố\rq\rq \ là
	 \choice
		{$2018$ không  chia hết cho $9$}
		{$2018$ không chia hết cho $18$} 
		{$2018$ không phải  là hợp số}
		{\True$2018$  không là số nguyên tố}
	 \loigiai{
		Phủ định của mệnh đề là \lq\lq $2018$  không là số nguyên tố\rq\rq.
	} 
\end{ex}

% nguồn: MaHoaNH_TracNghiem_Lop10
\begin{ex}%[0D1H1-3]%[ma:0524bf1d]
	Mệnh đề phủ định của mệnh đề  \lq\lq $2018$ là số tự nhiên chẵn\rq\rq\ là
	 \choice
		{$2018$  là số chẵn}
		{$2018$ là số nguyên  tố}
		{\True$2018$ không  là số tự nhiên chẵn}
		{$2018$ là số chính phương} 
	\loigiai{ 
		Mệnh đề phủ định là \lq\lq $2018$ không là  số tự nhiên chẵn\rq\rq.
	} 
\end{ex}

% nguồn: MaHoaNH_TracNghiem_Lop10
\begin{ex}%[0D1H1-3]%[ma:4a509110]
	Cho mệnh đề \lq\lq Có một học sinh trong lớp $10A$ không thích học môn Toán\rq\rq. Mệnh đề phủ định của  mệnh đề này là
	\choice 
		{\True\lq\lq Mọi học sinh trong lớp $10A$ đều thích  học môn Toán\rq\rq}
		{\lq\lq Mọi học sinh trong lớp $10A$ đều không thích học môn  Toán\rq\rq}
		 {\lq\lq Mọi học sinh trong lớp $10A$ đều thích học môn Văn\rq\rq}
		{\lq\lq  Có một học sinh trong lớp $10A$ thích học môn Toán\rq\rq}
	\loigiai{ 
		Mệnh đề phủ định là \lq\lq Mọi học sinh trong lớp $10A$  đều thích học môn Toán\rq\rq.
	 }
\end{ex}

% nguồn: MaHoaNH_TracNghiem_Lop10
\begin{ex}%[0D1H1-3]%[ma:6dd8d56d]
	Cho  mệnh đề \lq\lq Có một học sinh trong lớp $C4$ không chấp hành luật giao thông\rq\rq. Mệnh đề phủ định của mệnh đề này là
	\choice 
		{Không có học sinh nào  trong lớp $C4$ chấp hành luật giao thông}
		{\True Mọi học  sinh trong lớp $C4$ đều chấp hành luật giao thông}
		{Có một  học sinh trong lớp $C4$ chấp hành luật giao thông}
		{Mọi học sinh trong lớp $C4$ không  chấp hành luật giao thông}
	 \loigiai{
		Mệnh đề phủ định là \lq\lq Mọi học sinh trong lớp $C4$ đều chấp hành luật giao  thông\rq\rq.
	} 
\end{ex}

% nguồn: MaHoaNH_TracNghiem_Lop10
\begin{ex}%[0D1H1-3]%[ma:e87fee39]
	Cho mệnh đề $A$: \lq\lq $2$ là số nguyên  tố\rq\rq. Mệnh đề phủ định của mệnh đề $A$ là
	\choice 
		{$2$ không phải là  số hữu tỷ}
		 {$2$ là số nguyên}
		{\True$2$ không phải là số nguyên  tố}
		{$2$  là hợp số}
	\loigiai{ 
		$2$  không phải là số nguyên tố.
	} 
\end{ex}

% nguồn: MaHoaNH_TracNghiem_Lop10
\begin{ex}%[0D1H1-3]%[ma:8d0170c5]
	Phủ  định của mệnh đề \lq\lq $n > 9$\rq\rq\ là
	\choice 
		{\lq\lq  $-n > 9$\rq\rq}
		{\lq\lq $-n >-9$\rq\rq} 
		{\lq\lq $n < 9$  \rq\rq}
		{\True\lq\lq $n \leq 9$\rq\rq} 
	\loigiai{ 
		Phủ định của mệnh đề \lq\lq $n > 9$\rq\rq\  là \lq\lq $n \leq 9$\rq\rq.
	 }
\end{ex}

% nguồn: MaHoaNH_TracNghiem_Lop10
\begin{ex}%[0D1H1-3]%[ma:967e7fe1]
	Cho mệnh đề: \lq\lq  Có một học sinh trong lớp 10B không thích học môn Toán \rq\rq. Mệnh đề phủ định của mệnh đề này là
	\choice 
		{\lq\lq  Có một học sinh trong lớp 10B không thích học môn Toán \rq\rq}
		 {\True \lq\lq Mọi học sinh trong lớp 10B đều thích học môn Toán \rq\rq}
		{\lq\lq Mọi học sinh  trong lớp 10B đều không thích học môn Toán \rq\rq}
		{\lq\lq Có ít nhất một học sinh trong lớp 10B thích  học môn Toán \rq\rq}
	 \loigiai{
		Mệnh đề phủ định của mệnh đề  này là \lq\lq Mọi học sinh trong lớp 10B đều thích học môn Toán \rq\rq.
	 }
\end{ex}

% nguồn: MaHoaNH_TracNghiem_Lop10
\begin{ex}%[0D1H1-3]%[ma:4a734660]
	Cho mệnh đề \lq\lq Phương trình $x^2-4x+4=0$ có nghiệm\rq\rq. Mệnh đề  phủ định của mệnh đề đã cho và tính đúng, sai của mệnh đề phủ định là
	\choice 
		 {Phương trình $x^2-4x+4=0$ có nghiệm. Đây là mệnh đề đúng}
		{Phương  trình $x^2-4x+4=0$ có nghiệm. Đây là mệnh đề sai}
		{Phương  trình $x^2-4x+4=0$ vô nghiệm. Đây là mệnh đề đúng}
		{\True Phương trình $x^2-4x+4=0$ vô nghiệm. Đây là mệnh đề sai} 
	\loigiai{ 
		Mệnh đề phủ định  là \lq\lq Phương trình $x^2-4x+4=0$ vô nghiệm\rq\rq.\\
		Đây là mệnh đề sai vì  $x=2$ là nghiệm của phương trình.
	 }
\end{ex}

% nguồn: MaHoaNH_TracNghiem_Lop10
\begin{ex}%[0D1H1-3]%[ma:ed40a446]
	Mệnh đề nào sau  đây là mệnh đề phủ định của mệnh đề \lq\lq Mọi động vật đều di chuyển\rq\rq\ ?
	 \choice
		{Có ít nhất một  động vật di chuyển}
		{\True  Có ít nhất một động vật không di chuyển}
		 {Mọi động vật đều không di chuyển}
		{Mọi  động vật đều đứng yên}
	 \loigiai{
		Phủ định của \lq\lq mọi\rq\rq\ là  \lq\lq có ít nhất\rq\rq.\\
		Phủ định của \lq\lq đều  di chuyển\rq\rq\ là \lq\lq không di chuyển\rq\rq.\\
		Do đó mệnh đề phủ định của mệnh  đề \lq\lq Mọi động vật đều di chuyển\rq\rq\ là \lq\lq Có ít nhất một động vật không di chuyển\rq\rq.
	} 
\end{ex}

% nguồn: MaHoaNH_TracNghiem_Lop10
\begin{ex}%[0D1K1-3]%[ma:c21b2f78]
	Mệnh đề  phủ định của mệnh đề "$P: \forall x\in\mathbb{R}: x^2\geq 0$" là:
	 \choice{$\overline{P}$: "$\exists x\in\mathbb{R}: x^2\leq 0$"}
		{$\overline{P}$:  "$\forall x\in\mathbb{R}: x^2\leq 0$"}
		{\True $\overline{P}$:  "$\exists x\in\mathbb{R}: x^2<0$"}
		{$\overline{P}$: "$\forall x\notin\mathbb{R}: x^2\geq  0$"}
	\loigiai{Mệnh đề  phủ định của mệnh đề "$\forall x\in\mathbb{R}: x^2\geq 0$" là "$\exists x\in\mathbb{R}: x^2<0$".}
\end{ex}

% nguồn: MaHoaNH_TracNghiem_Lop10
\begin{ex}%[0D1K1-3]%[ma:b8aeca01]
	Mệnh đề phủ định của  mệnh đề $Q: \forall x\in\mathbb{R}: x^2+1\neq 0$" là:
	 \choice{\True $\overline{Q}$: "$\exists x\in\mathbb{R}: x^2+1=0$"}
		{$\overline{Q}$: "$\forall  x\notin\mathbb{R}: x^2+1\neq 0$"}
		 {$\overline{Q}$: "$\forall x\in\mathbb{R}: x^2+1=0$"}
		{$\overline{Q}$: "$\exists x\in\mathbb{R}: x^2+1\neq 0$"} 
	\loigiai{Mệnh đề phủ định  của mệnh đề "$\forall x\in\mathbb{R}: x^2+1\neq 0$" là "$\exists x\in\mathbb{R}: x^2+1=0$".}
\end{ex}

% nguồn: MaHoaNH_TracNghiem_Lop10
\begin{ex}%[0D1K1-3]%[ma:50a2e458]
	 Mệnh đề phủ định của mệnh đề "$P: \exists x\in\mathbb{R}: x-3>0$" là:
	\choice{$\overline{P}$: "$\exists x\in\mathbb{R}:  x-3\leq 0$"}
		{\True $\overline{P}$: "$\forall x\in\mathbb{R}: x-3\leq 0$"} 
		 {$\overline{P}$: "$\forall x\in\mathbb{R}: x-3>0$"}
		{$\overline{P}$: "$\exists x\notin\mathbb{R}:  x-3>0$"}
	\loigiai{Mệnh đề phủ định của mệnh đề "$\exists  x\in\mathbb{R}: x-3>0$" là "$\forall x\in\mathbb{R}: x-3\leq 0$".}
\end{ex}

% nguồn: MaHoaNH_TracNghiem_Lop10
\begin{ex}%[0D1B1-3]%[ma:f1f68625]
	Phủ định của mệnh đề ``Có ít nhất một số vô tỷ là số thập phân vô hạn tuần hoàn'' là mệnh đề  nào sau đây?
	\choice 
		{Mọi số vô tỷ đều là số thập  phân vô hạn tuần hoàn}
		{Có ít nhất một số vô tỷ là số thập phân vô hạn không  tuần hoàn}
		{\True Mọi số vô tỷ đều là số thập phân vô hạn không tuần hoàn} 
		{Mọi số vô tỷ  đều là số thập phân tuần hoàn}
	\loigiai{Phủ định của mệnh đề ``Có ít nhất một số vô tỷ là số thập  phân vô hạn tuần hoàn'' là mệnh đề ``Mọi số vô tỷ đều là số thập phân vô hạn không tuần hoàn''.}
\end{ex}

% nguồn: MaHoaNH_TracNghiem_Lop10
\begin{ex}%[0D1K1-3]%[ma:04c63a1f]
	Giả thuyết Goldbach khẳng định rằng mọi số nguyên chẵn lớn hơn $2$ đều có thể viết được dưới dạng tổng của hai số nguyên tố (chẳng hạn $2016=13+2003$). Và cho đến bây giờ, chưa có ai chứng minh được giả thuyết trên là  đúng, và cũng chưa có ai tìm được một phản ví dụ chỉ ra rằng giả thuyết trên là sai. Hỏi một phản ví dụ chứa nội dung nào dưới đây?
	\choice 
		{Một số nguyên lẻ lớn hơn $2$ mà có thể viết được  dưới dạng tổng của hai số nguyên tố}
		{Một số nguyên lẻ lớn hơn $2$  mà không thể viết được dưới dạng tổng của hai số nguyên tố}
		{Một số nguyên chẵn lớn hơn $2$ mà có thể viết được dưới dạng tổng của hai  số không nguyên tố}
		{\True  Một số nguyên chẵn lớn hơn $2$ mà không thể viết được dưới dạng tổng của hai số nguyên tố}
	\loigiai{Phản ví dụ là ví dụ chứa nội dung  ``Một số nguyên chẵn lớn hơn $2$ mà không thể viết được dưới dạng tổng của hai số nguyên tố''.}
\end{ex}

% nguồn: MaHoaNH_TracNghiem_Lop10
\begin{ex}%[0D1K1-3]%[ma:0ca71df3]
	Phát biểu mệnh đề phủ định của mệnh đề "Trong  vũ trụ tồn tại một hành tinh mà mọi địa điểm trên bề mặt hành tinh đó có nhiệt độ nhỏ hơn $-100^\circ $C".
	 \choice
		{"Trong vũ trụ tồn  tại một hành tinh mà mọi địa điểm trên bề mặt hành tinh đó có nhiệt độ lớn hơn $-100^\circ $C"}
		{"Trong  vũ trụ tồn tại một hành tinh mà mọi địa điểm trên bề mặt hành tinh đó có nhiệt độ không nhỏ hơn $-100^\circ $C"}
		{"Trong vũ trụ tồn tại một hành tinh có ít nhất một địa điểm trên bề mặt có nhiệt độ lớn  hơn hoặc bằng $-100^\circ $C"}
		{\True "Trong vũ trụ mọi hành tinh đều có ít nhất một địa điểm  trên bề mặt có nhiệt độ lớn hơn hoặc bằng $-100^\circ $C"}
	\loigiai{Mệnh đề phủ định của mệnh đề "Trong vũ trụ tồn tại một hành tinh mà mọi địa  điểm trên bề mặt hành tinh đó có nhiệt độ nhỏ hơn $-100^\circ $C" là "Trong vũ trụ mọi hành tinh đều có ít nhất một địa điểm trên bề mặt có nhiệt độ lớn hơn hoặc bằng $-100^\circ $C".}
\end{ex}

% nguồn: MaHoaNH_TracNghiem_Lop10
\begin{ex}%[0D1K1-3]%[ma:a990ffed]
	Phát biểu mệnh đề  phủ định của mệnh đề "Tồn tại một loài thú có nhiệt độ thân thể nhỏ hơn $35^\circ$C".
	\choice 
		{"Mọi loài thú có nhiệt độ thân thể  nhỏ hơn $35^\circ$C"}
		{"Tồn tại một loài thú có nhiệt  độ thân thể lớn hơn $35^\circ$C"}
		{\True "Mọi loài thú có nhiệt độ thân thể  không nhỏ hơn $35^\circ$C"}
		{"Mọi loài thú có nhiệt độ thân thể lớn hơn $35^\circ$C"} 
	\loigiai{Mệnh đề phủ định của mệnh đề "Tồn tại một loài thú có nhiệt độ thân thể nhỏ hơn $35^\circ$C" là  mệnh đề "Mọi loài thú có nhiệt độ thân thể không nhỏ hơn $35^\circ$C".}
\end{ex}

% nguồn: MaHoaNH_TracNghiem_Lop10
\begin{ex}%[0D1B1-3]%[ma:3a537372]
	Mệnh đề: “Mọi động vật đều di chuyển” có mệnh  đề phủ định là
	 \choice
		{Có  ít nhất một động vật di chuyển}
		{Mọi động  vật đều đứng yên}
		{\True Có ít nhất  một động vật không di chuyển}
		{Mọi động  vật đều không di chuyển}
	\loigiai{Phủ  định của mệnh đề trên là ``
	Có ít nhất  một động vật không di chuyển''.}
\end{ex}

% nguồn: MaHoaNH_TracNghiem_Lop10
\begin{ex}%[0D1Y1-3]%[ma:d95f6c06]
	 Cho mệnh đề $P\colon$ \lq\lq $9$ là số chia hết cho $3$\rq\rq. Mệnh đề phủ định của mệnh đề $P$ là
	\choice 
		{$\overline{P}\colon$\lq\lq  $9$ là bội của $3$\rq\rq}
		{\True $\overline{P}\colon$ \lq\lq  $9$ là số không chia hết cho $3$\rq\rq}
		{$\overline{P}\colon$\lq\lq $9$  là ước của $3$\rq\rq}
		{$\overline{P}\colon$\lq\lq $9$ là số lớn hơn $3$\rq\rq} 
	 \loigiai{
	Mệnh đề $P\colon$\lq\lq $9$ là số  chia hết cho $3$\rq\rq có mệnh đề phủ định là $\overline{P}\colon$ \lq\lq $9$ là số không chia hết cho $3$\rq\rq.}
\end{ex}

% nguồn: MaHoaNH_TracNghiem_Lop10
\begin{ex}%[0D1Y1-3]%[ma:04462893]
	Tìm  mệnh đề phủ định của mệnh đề sau \lq \lq $\forall x \in \mathbb{R},x^2 \geq 0$ \rq \rq.
	 \choice
		{\lq \lq $\forall  x \in \mathbb{R},x^2<0$ \rq \rq}
		 {\lq \lq $\forall x \in \mathbb{R},x^2>0$ \rq \rq}
		{\True \lq \lq $\exists x \in  \mathbb{R},x^2<0$ \rq \rq}
		{\lq \lq $\exists x \in \mathbb{R},x^2 \geq 0$ \rq  \rq}
	 \loigiai{
	Mệnh đề phủ định cần tìm là  \lq \lq $\exists x \in \mathbb{R},x^2<0$ \rq \rq.}
\end{ex}

% nguồn: MaHoaNH_TracNghiem_Lop10
\begin{ex}%[0D1B1-3]%[ma:20e12e16]
	 Cho mệnh đề "Phương trình $x^2+2x+1=0$ có nghiệm". Tìm mệnh đề phủ định của mệnh đề trên và cho biết tính đúng sai của mệnh đề phủ định.
	\choice 
		{``Phương trình $x^2+2x+1=0$ có hai nghiệm phân  biệt''. Đây là mệnh đề sai}
		{``Phương trình $x^2+2x+1=0$ có  nghiệm kép''. Đây là mệnh đề đúng}
		{\True  ``Phương trình $x^2+2x+1=0$ vô nghiệm''. Đây là mệnh đề sai}
		{``Phương trình $x^2+2x+1=0$ vô nghiệm''. Đây  là mệnh đề đúng}
	\loigiai{Mệnh đề phủ định là ``Phương trình $x^2+2x+1=0$  vô nghiệm.''\\
	Rõ ràng phương trình này có nghiệm $x=-1$ nên  mệnh đề này sai.}
\end{ex}

% nguồn: MaHoaNH_TracNghiem_Lop10
\begin{ex}%[0D1Y1-3]%[ma:5ad5f1d4]
	Cho mệnh đề \lq\lq$\forall x\in \mathbb{R},\, x^2+1>0$\rq \rq.  Mệnh đề phủ định của mệnh đề đã cho là
	 \choice
		{\True  \lq\lq$\exists x\in \mathbb{R},\, x^2+1\leq 0$\rq \rq}
		{\lq\lq$\exists x\in \mathbb{R},\,  x^2+1>0$\rq \rq}
		{\lq\lq$\forall x\in \mathbb{R},\, x^2+1\leq 0$\rq  \rq}
		 {\lq\lq$\forall x\in \mathbb{R},\, x^2+1<0$\rq \rq}
	 \loigiai{
	Mệnh đề phủ định của mệnh đề \lq\lq$\forall x\in \mathbb{R},\, x^2+1>0$\rq \rq\, là \lq\lq$\exists  x\in \mathbb{R},\, x^2+1\leq 0$\rq \rq.}
\end{ex}

% nguồn: MaHoaNH_TracNghiem_Lop10
\begin{ex}%[0D1B1-3]%[ma:1d835e61]
	Phủ định của mệnh đề $``\exists x\in \mathbb{Q}\colon 2x^2-5x+2=0"$  là
	 \choice
		{$``\exists  x\in \mathbb{Q}\colon 2x^2-5x+2>0"$}
		{$``\exists x\in \mathbb{Q}\colon 2x^2-5x+2\ne 0"$} 
		{\True $``\forall x\in  \mathbb{Q}\colon 2x^2-5x+2\ne 0"$}
		{$``\forall x\in \mathbb{Q}\colon  2x^2-5x+2=0"$}
	 \loigiai{
	Vì phủ định của mệnh đề  $``\exists x\in \mathbb{Q}\colon 2x^2-5x+2=0"$ là $``\forall x\in \mathbb{Q}\colon 2x^2-5x+2\ne 0"$.}
\end{ex}

% nguồn: MaHoaNH_TracNghiem_Lop10
\begin{ex}%[0D1B1-3]%[ma:72bf33cf]
	Tìm mệnh đề phủ định của  mệnh đề $P\colon$ \lq\lq$\forall x\in \mathbb{N}:x^2+x-1>0$\rq\rq.
	\choice 
		 {$\overline{P}\colon ``\forall x\in \mathbb{N}:x^2+x-1>0$''}
		{\True $\overline{P}\colon ``\exists x\in  \mathbb{N}:x^2+x-1\leqslant 0$''}
		 {$\overline{P}\colon ``\exists x\in \mathbb{N}:x^2+x-1>0$''}
		{$\overline{P}\colon ``\forall  x\in \mathbb{N}:x^2+x-1\leqslant 0$''}
	 \loigiai{
		Ta có $\overline{P}\colon$ ``$\exists x\in \mathbb{N}:x^2+x-1\le  0$''.
	 }
\end{ex}

% nguồn: MaHoaNH_TracNghiem_Lop10
\begin{ex}%[0D1Y1-3]%[ma:aed6bf7a]
	Giá trị gần đúng của $\sqrt{5}$ chính  xác đến hàng phần trăm là
	 \choice
		{$2{,}2$ } 
		{\True  $2{,}24$}
		 {$2{,}3$}
		{$2{,}236$} 
	 \loigiai{
		Ta có $\sqrt{5}\approx  2{,}24$.
	} 
\end{ex}

% nguồn: MaHoaNH_TracNghiem_Lop10
\begin{ex}%[0D1B1-3]%[ma:6a5366a3]
	Tìm mệnh đề phủ  định của mệnh đề: \lq\lq Dơi là một loài chim\rq\rq.
	\choice 
		{Dơi là một loài  thú có cánh}
		{\True  Dơi không phải là một loài chim}
		 {Chim cùng loài với dơi}
		 {Dơi là loài ăn sâu bọ}
	\loigiai{ 
		Để phủ định một mệnh đề, ta thêm (hoặc bớt) từ \lq\lq không \rq\rq (hoặc \lq\lq không phải\rq\rq)  vào trước vị ngữ của mệnh đề đó. Vậy mệnh đề phủ định của mệnh đề \lq\lq Dơi là một loài chim\rq\rq\, là \lq\lq Dơi không phải là một loài chim\rq\rq.
	 }
\end{ex}

% nguồn: MaHoaNH_TracNghiem_Lop10
\begin{ex}%[0D1Y1-3]%[ma:37caf91c]
	Mệnh đề phủ định của mệnh đề \lq \lq $2018$ là số tự nhiên chẵn\rq \rq \, là 
	\choice 
		{$  2018 $ là số chẵn}
		{\True $ 2018 $ không là  số tự nhiên chẵn}
		{$ 2018  $ là số nguyên tố}
		{$ 2018 $ là số chính  phương}
	 \loigiai{
		Phủ định của mệnh đề đã  cho là \lq \lq $2018$ không là số tự nhiên chẵn\rq \rq .
	} 
\end{ex}

% nguồn: MaHoaNH_TracNghiem_Lop10
\begin{ex}%[0D1B1-3]%[ma:fa6b8157]
	Mệnh đề nào dưới đây là  mệnh đề phủ định của mệnh đề “Hôm nay, trời nắng to”?
	 \choice
		{Hôm  qua, trời nắng to}
		 {Hôm nay, trời nắng không to}
		{\True Hôm nay,  trời không nắng to}
		 {Hôm nay, trời mưa to}
	\loigiai{ 
		Theo  định nghĩa, ta có mệnh đề phủ định của mệnh đề “Hôm nay, trời nắng to” là mệnh đề “Hôm nay, trời không nắng to”.
	 }
\end{ex}

% nguồn: MaHoaNH_TracNghiem_Lop10
\begin{ex}%[0D1B1-3]%[ma:61721137]
	 Cho mệnh đề $P(x)$: \lq\lq $\forall x \in \mathbb{R} : x^{2}+x+1>0$\rq\rq. Phủ định của mệnh đề $P(x)$ là
	\choice 
		{\True \lq\lq$\exists x  \in \mathbb{R} : x^{2}+x+1 \leq 0$\rq\rq}
		{\lq\lq$\forall x \in  \mathbb{R} : x^{2}+x+1<0$\rq\rq}
		 {\lq\lq$\forall x \in \mathbb{R} : x^{2}+x+1 \leq 0$\rq\rq}
		{\lq\lq$\exists x \in  \mathbb{R} : x^{2}+x+1<0$\rq\rq}
	 \loigiai{
		Mệnh đề phủ định $\overline{P(x)}$:  \lq\lq$\exists x \in \mathbb{R} : x^{2}+x+1 \leq 0$\rq\rq.
	 }
\end{ex}

% nguồn: MaHoaNH_TracNghiem_Lop10
\begin{ex}%[0D1B1-3]%[ma:5c0c30c1]
	Mệnh đề nào sau đây là phủ định của  mệnh đề “Mọi động vật đều di chuyển” ?
	 \choice
		{ Mọi động vật đều không di  chuyển}
		{ Mọi động vật đều đứng  yên}
		{ \True Có  ít nhất một động vật không di chuyển}
		{ Có ít nhất một động vật di  chuyển}
	 \loigiai{
		Theo định nghĩa, phủ định của mệnh đề “Mọi động vật đều di chuyển” là mệnh đề “Có ít nhất một động vật không di  chuyển”.
	 }
\end{ex}

% nguồn: MaHoaNH_TracNghiem_Lop10
\begin{ex}%[0D1B1-3]%[ma:b07e51b4]
	Cho mệnh đề $A\colon ``\forall x\in \Bbb{R}\colon x^2>x"$. Mệnh đề  phủ định của mệnh đề $A$ là
	\choice 
		{$\forall  x\in \Bbb{R}\colon x^2\le x$}
		{$\forall  x\in \Bbb{R}\colon x^2< x$}
		{\True $\exists x\in \Bbb{R}\colon x^2\le x$} 
		{$\exists\in \Bbb{R}\colon x^2\neq x$} 
	 \loigiai{
		Mệnh đề phủ định của mệnh đề $A$ là $\exists  x\in \Bbb{R}\colon x^2\le x$.
	} 
\end{ex}

% nguồn: MaHoaNH_TracNghiem_Lop10
\begin{ex}%[0D1B1-3]%[ma:3b39b8a7]
	Phủ định của mệnh đề “Dơi là một loài chim” là  mệnh đề nào sau đây?
	\choice 
		{Dơi  là một loài có cánh}
		{Chim cùng loài với  dơi}
		{Bồ câu là  một loài chim}
		{\True Dơi không  phải là một loài chim}
	 \loigiai{
		Theo định nghĩa, ta có mệnh đề phủ định của mệnh đề “Dơi là một loài chim” là mệnh đề “Dơi không phải là một loài  chim”.
	} 
\end{ex}

% nguồn: MaHoaNH_TracNghiem_Lop10
\begin{ex}%[0D1B1-3]%[ma:2341ec21]
	Mệnh đề phủ định của mệnh đề: \lq\lq $\exists x \in \mathbb{R}  , 2 x^2+3 x-5<0 $ \rq\rq \, là mệnh đề nào dưới đây?
	\choice 
		{$\forall x  \in \mathbb{R} , 2 x^2+3 x-5>0$}
		{\True \lq\lq $ \forall x \in \mathbb{R} , 2 x^2+3x-5  \geq 0 $ \rq\rq}
		{\lq\lq $\exists x \in \mathbb{R} | 2 x^2+3 x-5>0 $  \rq\rq}
		{\lq\lq $\exists x \in \mathbb{R} | 2 x^{2}+3 x-5  \geq 0$ \rq\rq}
	 \loigiai{
		Mệnh đề phủ định của mệnh đề \lq\lq$ \exists x \in \mathbb{R} , 2 x^2-3 x-5<0 $\rq\rq \, là mệnh đề là \lq\lq $\forall x \in \mathbb{R} , 2  x^2+3x-5 \geq 0 $\rq\rq.
	} 
\end{ex}

% nguồn: MaHoaNH_TracNghiem_Lop10
\begin{ex}%[0D1B1-3]%[ma:f913b3be]
	 Mệnh đề phủ định của mệnh đề \lq\lq $\exists x \in \mathbb{Z}, x \le \dfrac{1}{x}$ \rq \rq\ là
	\choice 
		{\lq \lq  $\forall x \in \mathbb{Z}, x \ge \dfrac{1}{x} $ \rq \rq}
		{\lq \lq $\exists x \in  \mathbb{Z}, x > \dfrac{1}{x}$ \rq \rq}
		{\True \lq  \lq $\forall x \in \mathbb{Z}, x >\dfrac{1}{x}$ \rq \rq}
		{\lq \lq $\exists  x \in \mathbb{Z}, x \le \dfrac{1}{x}$ \rq \rq}
	\loigiai{ 
		Phủ định của  mệnh đề \lq \lq $\exists x \in \mathbb{R}$ \rq \rq\ là \lq \lq $\forall x \in \mathbb{R}$ \rq \rq.\\
		Phủ định của \lq \lq $x \le  \dfrac{1}{x}$\rq \rq\ là \lq \lq $x > \dfrac{1}{x}$ \rq\rq.
	} 
\end{ex}

% nguồn: MaHoaNH_TracNghiem_Lop10
\begin{ex}%[0D1B1-3]%[ma:6ece21eb]
	Mệnh đề phủ định của mệnh đề $P \colon$  \lq\lq $\forall x \in \mathbb{R}, 3x^2+2>0$\rq\rq \text{ là}
	 \choice
		{\True $\overline{P} \colon$ \lq\lq  $\exists x \in \mathbb{R}, 3x^2+2 \le 0$\rq\rq}
		{$\overline{P} \colon$ \lq\lq $\forall x \in \mathbb{R}, 3x^2+2 \le  0$\rq\rq}
		{$\overline{P}  \colon$ \lq\lq $\exists x \in \mathbb{R}, 3x^2+2 < 0$\rq\rq}
		{$\overline{P} \colon$ \lq\lq $\exists x \in  \mathbb{R}, 3x^2+2 \ne 0$\rq\rq}
	\loigiai{ 
		Phủ định của \lq\lq $\forall x \in \mathbb{R}$\rq\rq \text{ là} \lq\lq $\exists x  \in \mathbb{R}$\rq\rq.\\
		Phủ định của \lq\lq $3x^2+2>0$\rq\rq  \text{ là} \lq\lq $3x^2+2 \le 0$\rq\rq.\\
		Do  vậy $\overline{P} \colon$ \lq\lq $\exists x \in \mathbb{R}, 3x^2+2 \le 0$\rq\rq.
	} 
\end{ex}

% nguồn: MaHoaNH_TracNghiem_Lop10
\begin{ex}%[0D1B1-3]%[ma:813b29d8]
	Trong các mệnh đề  sau, mệnh đề nào có mệnh đề phủ định \textbf{sai}?
	\choice 
		 {Phương trình $x^2-3x+2=0$ vô nghiệm}
		 {$2^{10}-1$ không chia hết cho $ 11$}
		{Có hữu  hạn số nguyên tố}
		{\True $7^{2017}-2^{2017}$  chia hết cho $5$}
	 \loigiai{
		Mệnh đề phủ định của  mệnh đề $P$ chỉ sai khi mệnh đề $P$ đúng.\\
		- Mệnh đề phương trình $x^2-3x+2=0$ vô nghiệm là  một mệnh đề sai.\\
		- Mệnh đề $2^{10}-1$ không chia hết cho $ 11$  là một mệnh đề sai vì $(2^{10}-1):11 =93$.\\
		- Mệnh đề có hữu hạn số nguyên tố  là một mệnh đề sai.\\
	- Mệnh đề  $7^{2017}-2^{2017}$ chia hết cho $5$ là một mệnh đề đúng vì $7^{2017}-2^{2017} = (7-2)(7^{2016}+7^{2015}.2+...+2^{2016})$.}
\end{ex}


%%%%%%%%%%%%%%%% ĐÚNG/SAI %%%%%%%%%%%%%%%%

% nguồn: MaHoaNH_DungSai_Lop10
\begin{ex}%[0D1N1-3]%[ma:39fe3110]
	Cho mệnh đề $P$: \lq\lq$23$ là số nguyên tố\rq\rq. Các câu  sau là đúng hay sai?
	 \choiceTF
		{\True Mệnh đề phủ định  của mệnh đề $P$ là $\overline{P}$: \lq\lq$23$ là hợp số\rq\rq}
		{\True Mệnh đề phủ định của mệnh đề $P$ là $\overline{P}$: \lq\lq$23$ không là  số nguyên tố\rq\rq}
		{Mệnh  đề phủ định của mệnh đề $P$ là mệnh đề đúng}
		{\True Mệnh đề  phủ định của mệnh đề $P$ là mệnh đề sai}
	 \loigiai{
	\begin{itemchoice}
		\itemch Mệnh đề phủ định của  mệnh đề $P$ là $\overline{P}$: \lq\lq$23$ là hợp số\rq\rq.
		\itemch  Mệnh đề phủ định của mệnh đề $P$ là $\overline{P}$: \lq\lq$23$ không là số nguyên tố\rq\rq.
		\itemch Mệnh đề  $P$ là mệnh đề đúng, do đó mệnh đề phủ định của mệnh đề $P$ là mệnh đề sai.
		\itemch Mệnh đề $P$ là mệnh đề đúng, do đó mệnh đề phủ định của mệnh đề $P$ là  mệnh đề sai.
	\end{itemchoice}
	 }
\end{ex}

% nguồn: MaHoaNH_DungSai_Lop10
\begin{ex}%[0D1H1-3]%[ma:3f5bc29e]
	Cho biết mệnh đề phủ định của mệnh đề sau  đúng hay sai?
	\choiceTF 
		{$P$: \lq\lq Hình thoi có hai đường chéo vuông góc với nhau\rq\rq. Ta  có mệnh đề phủ định là $\overline{P}$: \lq\lq Hình thoi có hai đường chéo không vuông góc với nhau\rq\rq}
		{$S$: \lq\lq $1>-3$\rq\rq.  Ta có mệnh đề phủ định là $\overline{S}$: \lq\lq $1\leq-3$\rq\rq}
		 {\True $K$: \lq\lq Phương trình $x^4-2x^2+2=0$ có nghiệm\rq\rq. Ta có mệnh đề phủ định là $\overline{K}$: \lq\lq Phương trình $x^4-2x^2+2=0$ vô nghiệm\rq\rq}
		{$H$: \lq\lq $(\sqrt{3}-\sqrt{12})^2=3$\rq\rq. Ta có mệnh đề phủ định là  $\overline{H}:$ \lq\lq $(\sqrt{3}-\sqrt{12})^2\neq 3$\rq\rq}
	\loigiai{ 
		\begin{itemchoice}
			 \itemch Mệnh đề này sai.
			 \itemch Mệnh đề này sai.
			\itemch Mệnh  đề này đúng
			\itemch Mệnh đề  này sai.
		\end{itemchoice}
	 }
\end{ex}

% nguồn: MaHoaNH_DungSai_Lop10
\begin{ex}%[0D1H1-3]%[ma:f37a1c7c]
	Xét tính đúng, sai của các câu  sau.
	\choiceTF 
		{\True  $P$: \lq\lq $3^3$ là số chính phương\rq\rq\ có mệnh đề phủ định là $\overline{P}$: \lq\lq $3^3$ không là số chính phương\rq\rq}
		 {$Q$: \lq\lq Tam giác $ABC$ là tam giác cân\rq\rq\ có mệnh đề phủ định là $\overline{Q}$: \lq\lq Tam giác $ABC$ không là tam giác vuông\rq\rq}
		{\True $R$: \lq\lq  $2^{2003}-1$ là số nguyên tố\rq\rq\ có mệnh đề phủ định là $\overline{R}$: \lq\lq $2^{2003}-1$ không là số nguyên tố\rq\rq}
		{\True $H$: \lq\lq $\sqrt{2}$ là số vô tỉ\rq\rq\ có mệnh đề phủ định là  $\overline{H}$: \lq\lq $\sqrt{2}$ là số hữu tỉ\rq\rq}
	\loigiai{ 
		\begin{itemchoice}
			\itemch $\overline{P}$: \lq\lq  $3^3$ không là số chính phương\rq\rq.
			\itemch $\overline{Q}$: \lq\lq Tam giác $ABC$  không là tam giác cân\rq\rq.
			\itemch $\overline{R}$: \lq\lq  $2^{2003}-1$ không là số nguyên tố\rq\rq.
			\itemch $\overline{H}$: \lq\lq $\sqrt{2}$ là số hữu tỉ\rq\rq. 
		\end{itemchoice}
	} 
\end{ex}

% nguồn: MaHoaNH_DungSai_Lop10
\begin{ex}%[0D1H1-3]%[ma:36cc6af3]
	Cho biết mệnh đề  phủ định của mệnh đề sau đúng hay sai?
	\choiceTF 
		{$A$: \lq\lq $-\dfrac{1{,}3}{5}$ là một phân số\rq\rq\ và  $\overline{A}$: \lq\lq $-\dfrac{1{,}3}{5}$ là số tự nhiên\rq\rq}
		{\True $B$:\lq\lq Phương trình $x^2+3x-2023=0$ có nghiệm\rq\rq\ và $\overline{B}$: \lq\lq Phương  trình $x^2+3x-2023=0$ không có nghiệm\rq\rq\ là mệnh đề sai}
		{\True $D$:\lq\lq Số $2023$ chia hết cho $17$\rq\rq\ và $\overline{D}$: \lq\lq Số $2023$ không chia  hết cho $17$\rq\rq}
		{$F$: \lq\lq Hai đường thẳng $y=2023x+1$ và $y=-2023x+1$  không song song với nhau\rq\rq\ và $\overline{F}$: \lq\lq Hai đường thẳng $y=2023x+1$ và $y=-2023x+1$ vuông góc với nhau\rq\rq}
	\loigiai{ 
		\begin{itemchoice}
			 \itemch Mệnh đề phủ định của mệnh đề $A$ là $\overline{A}$ \lq\lq $-\dfrac{1{,}3}{5}$ không là phân số\rq\rq.\\
			Mệnh đề $\overline{A}$ đúng  vì $-1{,}3$ không là số nguyên.
			\itemch Mệnh đề phủ định của  mệnh đề $B$ là $\overline{B}$: \lq\lq Phương trình $x^2+3x-2023=0$ không có nghiệm\rq\rq.\\
			Mệnh đề $\overline{B}$ sai vì phương  trình $x^2+3x-2023=0$ có hai nghiệm phân biệt do $a \cdot c=1\cdot(-2023) < 0$.
			\itemch Mệnh đề phủ định của mệnh đề $D$ là $\overline{D}$: \lq\lq Số $2023$  không chia hết cho $17$\rq\rq.\\
			Mệnh đề $\overline{D}$  sai vì $2023$ chia hết cho $17$.
			\itemch Mệnh đề phủ định của mệnh đề $F$ là $\overline{F}$: \lq\lq Hai đường  thẳng $y=2023x+1$ và $y=-2023x+1$ song song với nhau\rq\rq.\\ Mệnh đề $\overline{F}$ sai vì hai đường thẳng $(d_1)\colon y=2023x+1$ và $(d_2)\colon y=-2023x+1$ có hệ số góc $k_1 \neq k_2$ $(2023\neq-2023)$.
		\end{itemchoice}
	 }
\end{ex}

% nguồn: MaHoaNH_DungSai_Lop10
\begin{ex}%[0D1H1-3]%[ma:f76fc24c]
	Cho  mệnh đề $P(x):$ \lq\lq $\forall x\in \mathbb{R},x^2-4x+1> 0$\rq\rq.
	 \choiceTF
		{$P(x)$: \lq\lq Có ít nhất một số thực $x$ thoả mãn $x^2-4x+1> 0$\rq\rq} 
		{$P(x)$ là mệnh đề chứa biến} 
		{\True  $\overline{P(x)}$: \lq\lq $\exists x\in \mathbb{R}, x^2-4x+1\le 0$\rq\rq}
		{Mệnh đề $\overline{P(x)}$ là mệnh  đề sai}
	 \loigiai{
		\begin{itemchoice}
			\itemch  SAI.\\
			Ta có phát biểu đúng của mệnh đề $P(x):$\lq\lq Với mọi số thực  $x$ thoả mãn $x^2-4x+1> 0$\rq\rq.
			\itemch SAI.\\ 
			$P(x)$ là mệnh đề sai nên không phải là  mệnh đề chứa biến.
			\itemch  ĐÚNG.\\
			Mệnh đề phủ định  của $P(x)$ là $\overline{P(x)}:$ \lq\lq $\exists x\in \mathbb{R}, x^2-4x+1\le 0$\rq\rq.
			\itemch SAI.\\ 
			Ta có $\overline{P(x)}:$ \lq\lq $\exists x\in \mathbb{R}, x^2-4x+1\le 0$\rq\rq. Mệnh đề $\overline{P(x)}$ đúng vì giả sử ta  chọn $x=1$ thì $1^2-4.1+1=-2< 0$.
		\end{itemchoice}
	} 
\end{ex}


%%%%%%%%%%%%%%%% TRẢ LỜI NGẮN %%%%%%%%%%%%%%%%

% nguồn: MaHoaNH_TraLoiNgan_Lop10
\begin{bt}%[0D1H1-3]%[ma:b5f5e06e]
	Cho mệnh đề $P$: \lq\lq $x^2-3x+4=0$  vô nghiệm\rq\rq\ và các mệnh đề sau.
	\begin{itemize}
		\item  \lq\lq $x^2-3x+4=0$ có nghiệm\rq\rq.
		\item \lq\lq $x^2-3x+4=0$ có hai nghiệm phân  biệt\rq\rq.
		\item \lq\lq $x^2-3x+4=0$ không vô nghiệm\rq\rq. 
	\end{itemize}
	 Có bao nhiêu phát biểu là phủ định của mệnh đề $P$?
	\shortans{2} 
	\loigiai{ 
		Phủ định của mệnh đề $P$ là $\overline{P}$: \lq\lq $x^2-3x+4=0$ có nghiệm\rq\rq\ hoặc $\overline{P}$: \lq\lq $x^2-3x+4=0$  không vô nghiệm\rq\rq.\\
		Vậy có $2$ mệnh đề thỏa. Các mệnh đề còn lại không phải là phủ định của mệnh đề $P$. 
	 }
\end{bt}

% nguồn: MaHoaNH_TraLoiNgan_Lop10
\begin{bt}%[0D1H1-3]%[ma:d05bebba]
	Cho mệnh đề $P$: \lq\lq $x^2+x+1$ là số dương với $x$ là số thực\rq\rq\ và các  mệnh đề sau.
	\begin{itemize}
		\item \lq\lq $x^2+x+1$  là số không dương với $x$ là số thực\rq\rq.
		\item \lq\lq $x^2+x+1$ là khác $0$ với $x$ là số thực\rq\rq. 
		\item \lq\lq $x^2+x+1$ là số âm với $x$ là số thực\rq\rq. 
		\item \lq\lq $x^2+x+1$  là số không âm với $x$ là số thực\rq\rq.
	\end{itemize}
	Có bao nhiêu phát biểu là phủ định của mệnh  đề $P$?
	\shortans{1} 
	 \loigiai{
		Phủ định của mệnh đề $P$ là $\overline{P}$: \lq\lq $x^2+x+1$ là số không dương  với $x$ là số thực\rq\rq.\\
		Vậy  có $1$ mệnh đề thỏa. Các mệnh đề còn lại không phải là phủ định của mệnh đề $P$.
	 }
\end{bt}

% nguồn: MaHoaNH_TraLoiNgan_Lop10
\begin{bt}%[0D1H1-3]%[ma:6c4b8743]
	Cho mệnh đề $P$: \lq\lq $4$  là số chẵn\rq\rq\ và các mệnh đề sau.
	\begin{itemize}
		\item \lq\lq $4$ không là số  chẵn\rq\rq.
		\item \lq\lq $4$ là số  lẻ\rq\rq.
		\item \lq\lq $4$ không là số tự  nhiên lẻ\rq\rq.
		\item  \lq\lq $4$ không là số chia hết cho $2$\rq\rq.
	\end{itemize}
	Có bao nhiêu phát biểu là phủ định của mệnh đề  $P$?
	\shortans{2} 
	\loigiai{ 
		Phủ định của mệnh đề $P$ là $\overline{P}$: \lq\lq $4$ không là số chẵn\rq\rq\ hoặc $\overline{P}$: \lq\lq $4$  là số lẻ\rq\rq.\\
		Vậy có $2$ mệnh đề thỏa. Các mệnh đề còn lại không phải  là phủ định của mệnh đề $P$.
	 }
\end{bt}

% nguồn: MaHoaNH_TraLoiNgan_Lop10
\begin{bt}%[0D1K1-3]%[ma:7f419606]
	Cho mệnh đề $P\colon$ \lq\lq Bình phương của một số thực lớn hơn $ 2 $ \rq\rq, tìm số số  nguyên thỏa mãn mệnh đề phủ định của mệnh đề $P$.
	\shortans{3} 
	 \loigiai{
		Ta  có: $P\colon$ \lq\lq $x^2>2$\rq\rq.
		Suy ra $\overline{P}\colon$  \lq\lq $ x^2 \leq 2 $\rq\rq.\\
		Vì $x \in \mathbb{Z}$ nên $x$ là $ -1 $; $ 0 $;  $ 1 $.\\
	Vậy có $ 3 $ số nguyên  thỏa mãn mệnh đề $\overline{P}$.}
\end{bt}


%%%%%%%%%%%%%%%% TỰ LUẬN %%%%%%%%%%%%%%%%

% nguồn: MaHoaNH_TuLuan_Lop10
\begin{bt}%[0D1H1-3]%[ma:ac0b3892]
	Lập mệnh  đề phủ định của mỗi mệnh đề sau:
	\begin{enumerate}
		\item Mọi hình vuông đều là hình thoi. 
		\item Có một tam giác cân không phải là tam  giác đều.
	\end{enumerate}
	 \loigiai{
		Ta có các mệnh đề phủ định  là
		\begin{enumerate}
			\item  Có ít nhất một hình vuông không phải là hình thoi.
			\item Mọi tam giác cân đều là tam  giác đều.
		\end{enumerate}
	} 
\end{bt}

% nguồn: MaHoaNH_TuLuan_Lop10
\begin{bt}%[0D1H1-3]%[ma:669d404b]
	Lập  mệnh đề phủ định của mỗi mệnh đề sau:
	\begin{itemize}
		\item  $P$: \lq\lq Trong tam giác tổng ba góc bằng $180^{\circ}$ \rq\rq.
		\item $Q$:  \lq\lq $6$ không phải là số nguyên tố\rq\rq.
	\end{itemize}
	\loigiai{ 
		Ta có các mệnh đề phủ định  là
		\begin{itemize}
			\item $\overline{P}$: \lq\lq Trong tam giác  tổng ba góc không bằng $180^{\circ}$\rq\rq.
			\item $\overline{Q}:$ \lq\lq $6$ là số  nguyên tố\rq\rq.
		\end{itemize}
	 }
\end{bt}

% nguồn: MaHoaNH_TuLuan_Lop10
\begin{bt}%[0D1H1-3]%[ma:d4b469b4]
	Lập mệnh đề phủ định của mỗi mệnh đề sau và nhận  xét tính đúng sai của mệnh đề phủ định đó.
	\begin{enumerate}
		\item $A$: \lq\lq $\dfrac{5}{1{,}2}$  là một phân số\rq\rq.
		\item $B$: \lq\lq Phương  trình $x^2+3x+2=0$ có nghiệm\rq\rq.
		 \item $C$: \lq\lq $2^2+2^3=2^{2+3}$\rq\rq.
		\item $D$: \lq\lq Số $2025$ chia hết  cho $15$\rq\rq.
	\end{enumerate}
	 \loigiai{
		\begin{enumerate}
			\item  $\overline{A}$: \lq\lq$\dfrac{5}{1{,}2}$ không là một phân số\rq\rq.\\
			Mệnh đề đúng vì $\dfrac{5}{1{,}2}$ không là phân số (do $1{,}2$  không là số nguyên)
			\item $\overline{B}$: \lq\lq Phương trình  $x^2+3x+2=0$ vô nghiệm\rq\rq.\\
			Mệnh đề sai vì phương trình $x^2+3x+2=0$ có hai nghiệm là $x=-1$ và $x=-2$. 
			 \item $\overline{C}$: \lq\lq $2^2+2^3\neq 2^{2+3}$\rq\rq.\\
			Mệnh đề đúng vì $2^2+2^3=12\neq 32=2^{2+3}$. 
			\item $\overline{D}$: \lq\lq Số $2025$ không chia hết  cho $15$\rq\rq.\\
			Mệnh đề sai vì $2025$ chia hết cho $15$. 
		\end{enumerate}
	 }
\end{bt}

% nguồn: MaHoaNH_TuLuan_Lop10
\begin{bt}%[0D1H1-3]%[ma:287624eb]
	Lập mệnh đề phủ định của mỗi mệnh đề sau và nhận xét tính  đúng sai của mệnh đề phủ định đó.
	\begin{itemize}
		\item $A$: \lq\lq Hình thoi có hai  đường chéo vuông góc với nhau\rq\rq.
		\item $B$: \lq\lq $6$ là số nguyên  tố\rq\rq.
		\item $C$: \lq\lq Tổng hai cạnh của một tam giác lớn hơn  cạnh còn lại\rq\rq.
		\item $D$: \lq\lq  $5>-3$\rq\rq.
		\item  $E$: \lq\lq Phương trình $x^4-2x^2+2=0$ có nghiệm\rq\rq.
		\item $F$:  \lq\lq $\left(\sqrt{3}-\sqrt{12}\right)^2=3$\rq\rq.
	\end{itemize}
	\loigiai{ 
		 Ta có các mệnh đề phủ định là
		\begin{itemize}
			\item $\overline{A}:$ \lq\lq Hai đường chéo của hình thoi  không vuông góc với nhau\rq\rq, mệnh đề này sai.
			\item $\overline{B}$: \lq\lq $6$ không phải là  số nguyên tố\rq\rq, mệnh đề này đúng.
			\item $\overline{C}:$ \lq\lq Tổng hai cạnh của một tam giác nhỏ hơn hoặc bằng cạnh còn lại\rq\rq, mệnh  đề này sai.
			\item $\overline{D}$: \lq\lq $5\leq-3$ \rq\rq, mệnh đề  này sai.
			\item $\overline{E}$: \lq\lq Phương trình $x^4-2x^2+2=0$ vô  nghiệm\rq\rq, mệnh đề này đúng.
			\item $\overline{F}$: \lq\lq  $\left(\sqrt{3}-\sqrt{12}\right)^2\neq 3$\rq\rq, mệnh đề này sai.
		\end{itemize}
	} 
\end{bt}

% nguồn: MaHoaNH_TuLuan_Lop10
\begin{bt}%[0D1V1-3]%[ma:cf8fa871]
	Cho mệnh đề  $P: \forall n\in N: n > \dfrac{1}{n}$. Gọi $S$ là tập hợp các phần tử $n$ để mệnh đề $\overline{P}$ đúng. Tìm số phần tử của tập $S$.
	\loigiai{ 
		Ta có 
		\[\overline{P}: \exists  n\in N: n\le \dfrac{1}{n}
		\] 
		Vì $n\in N$ nên bất phương trình  xác định khi $n > 0$ hay $n\ge 1$. Ta có $n\le \dfrac{1}{n} \Leftrightarrow n^2\le 1\Leftrightarrow n\le 1$.\\
		Kết hợp  điều kiện ta có $1\le n\le 1\Rightarrow n=1$.\\
		 Vậy tập $S$ có đúng $1$ phần tử.
	} 
\end{bt}

% nguồn: MaHoaNH_TuLuan_Lop10
\begin{bt}%[0D1B1-3]%[ma:5b1951df]
	Nêu mệnh đề phủ định của các mệnh đề sau  và cho biết tính đúng sai của mệnh đề đó.
	\begin{center}
		$A: \lq\lq \exists x\in\mathbb{N}, n^2+3$ chia hết cho $4$ \rq\rq và $B: \lq\lq \exists x\in\mathbb{N}, x$ chia hết cho $x+1$  \rq\rq.
	\end{center}
	\loigiai{\begin{listEX}[1]
		\item $\overline{A}: \lq\lq \forall x\in\mathbb{N}, n^2+3$ không chia hết cho  $4$ \rq\rq - Sai.
		\item  $\overline{B}: \lq\lq \forall x\in\mathbb{N}, x$ không chia hết cho $x+1$ \rq\rq - Sai.
\end{listEX}}
\end{bt}

% nguồn: MaHoaNH_TuLuan_Lop10
\begin{bt}%[0D1B1-3]%[ma:b557d02b]
	Cho mệnh đề $P\colon$ \lq\lq$\exists  n\in \mathbb{Z}:6n^2-5n-4=0$\rq\rq.
	Xét tính \textbf{đúng - sai}, giải thích  và lập mệnh đề phủ định của mệnh đề $P$.
	\loigiai{ 
	Ta có: 
	$6n^2-5n-4=0\Leftrightarrow \hoac{& n=-\dfrac{1}{2}\notin \mathbb{Z}  \\ &n=\dfrac{4}{3}\notin \mathbb{Z} }
	\Rightarrow  P$ là mệnh đề \textbf{sai}.\\
	Mệnh đề phủ  định của mệnh đề $P$ là
	$\overline{P}\colon$ \lq\lq $\forall n\in \mathbb{Z}:6n^2-5n-4\neq 0  $\rq\rq.
	} 
\end{bt}

% nguồn: MaHoaNH_TuLuan_Lop10
\begin{bt}%[0D1B1-3]%[ma:cc03c9a6]
	Nêu mệnh đề phủ định của các mệnh đề sau và cho biết tính đúng sai của mệnh  đề phủ định đó.
	\begin{enumEX}{1}
		\item $A: \lq\lq \forall  x\in\mathbb{R}, x^3-x^2+1>0$ \rq\rq;
		\item $B:$ \lq\lq Tồn tại số thực $a$ sao cho $a+\dfrac 1 a\leq 2$ \rq\rq. 
	\end{enumEX}
	\loigiai{	\begin{enumEX}{1}
		\item  $\overline{A}: \lq\lq \exists x\in\mathbb{R}, x^3-x^2+1 \leq 0$ \rq\rq - Đúng;
		\item $\overline{B}:$ \lq\lq Với mọi số thực $a$ sao cho $a+\dfrac 1  a> 2$ \rq\rq - Sai.
\end{enumEX} }
\end{bt}

% nguồn: MaHoaNH_TuLuan_Lop10
\begin{bt}%[0D1B1-3]%[ma:55e3decc]
	Phát biểu mệnh đề phủ định  của các mệnh đề sau:
	\begin{enumerate}[a)]
		\item $P\colon$ \lq\lq Tháng $12$ dương lịch  có $31$ ngày\rq\rq;
		\item $Q\colon$ \lq\lq  $9^{10} \ge 10^9$\rq\rq;
		\item  $R\colon$ \lq\lq Phương trình $x^2+1=0$ có nghiệm\rq\rq.
	\end{enumerate}
	 \loigiai{
		\begin{enumerate}[a)]
			\item $\overline{P}\colon$ \lq\lq Không phải tháng $12$ dương lịch có  $31$ ngày\rq\rq.
			\item  $\overline{Q}\colon$ \lq\lq$9^{10}<10^9$\rq\rq.
			\item $\overline{R}\colon$ \lq\lq  Phương trình $x^2+1=0$ vô nghiệm\rq\rq.
		\end{enumerate}
	 }
\end{bt}

% nguồn: MaHoaNH_TuLuan_Lop10
\begin{bt}%[0D1B1-3]%[ma:4f5756b1]
	Nêu mệnh đề phủ định cúa các mệnh  đề sau và cho biết tính đúng sai của mệnh đề phủ định đó.
	\begin{enumerate}
		\item $\overline{A}:$ \lq\lq Phương trình  $x^4-2x^2+2=0$ có nghiệm\rq\rq;
		\item $\overline{B}:$  \lq\lq Bất phương trình $x^{2013}>2030$ vô nghiệm \rq\rq;
		\item  $\overline{C}:$ \lq\lq $\forall x\in\mathbb{R}, x^4-x^2+1=\left(x^2+\sqrt 3x+1\right)\left(x^2-\sqrt 3x+1\right)$ \rq\rq;
		\item $\overline{D}:$ \lq\lq $\exists q\in\mathbb{Q}, 2q^2-1=0  $ \rq\rq.
	\end{enumerate}
	\loigiai{\begin{enumerate}
		\item $A:$ \lq\lq Phương trình $x^4-2x^2+2=0$ vô  nghiệm\rq\rq - Đúng;
		\item  $B:$ \lq\lq Bất phương trình $x^{2013}> 2030$ có nghiệm \rq\rq- Đúng;
		\item $C:$ \lq\lq $\exists x\in\mathbb{R}, x^4-x^2+1\neq\left(x^2+\sqrt 3x+1\right)\left(x^2-\sqrt 3x+1\right)$ \rq\rq  - Sai;
		\item $D:$ \lq\lq $\forall q\in\mathbb{Q}, 2q^2-1\neq0 $ \rq\rq  - Đúng.
\end{enumerate}}
\end{bt}

% nguồn: MaHoaNH_TuLuan_Lop10
\begin{bt}%[0D1B1-3]%[ma:2e5a6bed]
	Nêu mệnh đề phủ định của các mệnh  đề sau và cho biết tính đúng sai của mệnh đề phủ định đó.
	\begin{enumEX}{1}
		\item $A:$ \lq\lq 6 là  số nguyên tố \rq\rq.
		\item $B:$ \lq\lq $(\sqrt 3-1)^2$  là số nguyên \rq\rq;
		\item $C:$ \lq\lq $\exists n\in\mathbb{N}, n(n+1)$ là số chính phương \rq\rq; 
		 \item $D:$ \lq\lq $\forall n\in\mathbb{N}, 2n+1$ là số lẻ \rq\rq.
	\end{enumEX}
	\loigiai{	\begin{enumEX}{1}
		\item $\overline{A}:$ \lq\lq 6  là hợp số\rq\rq - Đúng.
		\item  $\overline{B}:$ \lq\lq $(\sqrt 3-1)^2$ không phải là số nguyên \rq\rq - Đúng;
		\item $\overline{C}:$ \lq\lq $\forall  n\in\mathbb{N}, n(n+1)$ không phải là số chính phương \rq\rq - Sai;
		\item $\overline{D}:$ \lq\lq $\exists n\in\mathbb{N}, 2n+1$ là  số chẵn \rq\rq - Sai.
\end{enumEX}}
\end{bt}

% nguồn: MaHoaNH_TuLuan_Lop10
\begin{bt}%[0D1B1-3]%[ma:4dd50a7e]
	Xét tính đúng sai của các mệnh  đề sau và nêu mệnh đề phủ định của nó
	\begin{enumerate}
		\item $P(x): \lq\lq\exists x\in\mathbb{Z},  x^2=3$ \rq\rq.
		\item $P(n): \lq\lq \forall n\in\mathbb{N}^{*}:  2^n+3$ là một số nguyên tố \rq\rq.
		\item $P(x): \lq\lq  \forall x\in\mathbb{R}, x^2+4x+5>0$ \rq\rq.
		\item $P(x):  \lq\lq \forall x\in\mathbb{R}, x^4-x^2+2x+2\geq 0$ \rq\rq.
	\end{enumerate}
	\loigiai{\begin{enumerate}
		\item Sai và $\overline{P(x)}:  \lq\lq\forall x\in\mathbb{Z}, x^2\neq3$ \rq\rq.
		\item Sai và $\overline{P(n)}: \lq\lq \exists  n\in\mathbb{N}^{*}: 2^n+3$ là một hợp số \rq\rq.
		\item Đúng và $\overline{P(x)}: \lq\lq \exists  x\in\mathbb{R}, x^2+4x+5\leq 0$ \rq\rq.
		\item Đúng và $\overline{P(x)}: \lq\lq \exists  x\in\mathbb{R}, x^4-x^2+2x+2< 0$ \rq\rq.
\end{enumerate}}
\end{bt}

% nguồn: MaHoaNH_TuLuan_Lop10
\begin{bt}%[0D1B1-3]%[ma:0635268d]
	Nêu  mệnh đề phủ định của mỗi mệnh đề sau và xác định xem mệnh đề phủ định đó đúng hay sai.
	\begin{listEX}
		\item  Phương trình $x^2-3x+2=0$ có nghiệm.
		\item $2^{10}-1$ chia hết cho $11$. 
		\item Có  vô số số nguyên tố.
	\end{listEX}
	 \loigiai{
		\begin{listEX}
			\item Phương trình $x^2-3x+2=0$ có nghiệm. Mệnh đề phủ định  là phương trình $x^2-3x+2=0$ vô có nghiệm. Mệnh đề phủ đinh sai.
			\item $2^{10}-1$ chia hết cho $11$. Mệnh đề phủ định là $2^{10}-1$ không chia hết cho $11$. Mệnh đề phủ định sai. 
			\item Có vô số số nguyên tố. Mệnh đề phủ định là có hữu hạn hoặc không có số nguyên tố. Mệnh đề phủ định  sai.
		\end{listEX}
	} 
\end{bt}

% nguồn: MaHoaNH_TuLuan_Lop10
\begin{bt}%[0D1B1-3]%[ma:e5cbce1e]
	Nêu mệnh đề phủ định của các mệnh đề sau 
	\begin{enumEX}{1}
		\item Số tự  nhiên $n$ chia hết cho $2$ và cho $3$.
		\item Số tự nhiên $n$ có chữ số tận cùng bằng $0$ hoặc  bằng $5$.
		\item Tứ giác $T$ có hai cạnh đối vừa  song song vừa bằng nhau.
		\item Số tự nhiên $n$ có ước số bằng $1$ và bằng  $n$.
	\end{enumEX}
	 \loigiai{
		\begin{enumEX}{1}
			 \item Số tự nhiên $n$ không chia hết cho $2$ và cho $3$.
			\item Số tự nhiên $n$ không có chữ số tận cùng bằng $0$  hoặc bằng $5$.
			\item Tứ giác $T$ không có hai cạnh đối vừa song  song vừa bằng nhau.
			\item Số  tự nhiên $n$ không có ước số bằng $1$ và bằng $n$.
		\end{enumEX}
	} 
\end{bt}

% nguồn: MaHoaNH_TuLuan_Lop10
\begin{bt}%[0D1B1-3]%[ma:d6833cfc]
	Viết mệnh đề phủ định của mỗi mệnh đề sau và xét xem mệnh đề phủ định đó  đúng hay sai?
	\begin{listEX}[2]
		\item $\pi  <3{,}15$.
		 \item $|-125|\le 0$.
		\item $3$ là  số nguyên tố.
		\item $7$ không chia hết cho $5$. 
		\item $\pi$ là số hữu tỉ. 
		\item  $1794$ chia hết cho $3$.
		\item $\sqrt{2}$ là  số hữu tỉ.
		\item Tổng hai cạnh của một  tam giác lớn hơn cạnh thứ ba.
	\end{listEX}
	 \loigiai{
		\begin{listEX}
			\item $\pi <3{,}15$. Mệnh đề phủ định là $\pi \ge3{,}15$. Mệnh đề phủ  định sai.
			 \item $|-125|\le 0$. Mệnh đề phủ định là $|-125|> 0$. Mệnh đề phủ định đúng.
			\item $3$ là số nguyên tố. Mệnh đề phủ định là $3$ không là số nguyên  tố. Mệnh đề phủ định sai.
			\item $7$ không chia hết cho $5$. Mệnh đề phủ định là  $7$ chia hết cho $5$. Mệnh đề phủ định sai.
			\item $\pi$ là  số hữu tỉ. Mệnh đề phủ định là $\pi$ không là số hữu tỉ. Mệnh đề phủ định đúng.
			\item $1794$ chia hết cho $3$. Mệnh đề phủ định là $1794$ không chia hết cho $3$. Mệnh đề phủ định sai. 
			\item $\sqrt{2}$ là số hữu tỉ. Mệnh đề phủ định là $\sqrt{2}$ không là số hữu tỉ. Mệnh  đề phủ định đúng.
			\item Tổng hai cạnh của một tam giác lớn hơn cạnh thứ ba. Mệnh đề phủ định là tổng hai cạnh của một tam giác bé hơn hoặc  bằng cạnh thứ ba. Mệnh đề phủ định sai.
		\end{listEX}
	 }
\end{bt}
